\documentclass[aspectratio=169,11pt]{beamer}

\usetheme{Madrid}
\usecolortheme{default}
\usefonttheme{professionalfonts}
\setbeamertemplate{navigation symbols}{}
\setbeamertemplate{footline}[frame number]

\usepackage{amsmath, amssymb, booktabs}

\newcommand{\E}{\mathbb{E}}
\newcommand{\1}{\mathbf{1}}
\newcommand{\Var}{\operatorname{Var}}
\newcommand{\Cov}{\operatorname{Cov}}

\title{Causal Inference With Spatial Data}
\subtitle{Empirical Methods --- Week 16}
\author{Teaching deck draft}
\date{}

\begin{document}

\begin{frame}
  \titlepage
\end{frame}

\begin{frame}{Week 16 question}
\begin{itemize}
    \item How do spatial spillovers, sorting, exposure, and correlated shocks change causal designs?
    \item Lecture 15 built spatial objects; Lecture 16 asks when those objects support causal claims.
    \item The core discipline: name the geography, the counterfactual, the interference structure, and the inference problem.
    \item A map is not a design until the comparison and assumptions are visible.
\end{itemize}
\end{frame}

\begin{frame}{Geography has three jobs}
\small
\begin{tabular}{p{0.25\linewidth} p{0.34\linewidth} p{0.31\linewidth}}
\toprule
Role & What geography does & Paper template \\
\midrule
Source & border, distance, or exposure creates variation & Black; Dell; Dube--Lester--Reich \\
Threat & nearby units share residual shocks & Conley \\
Mechanism & access, neighborhood, spillover, or local labor market is the channel & Miguel--Kremer; Chetty et al. \\
\bottomrule
\end{tabular}

\medskip
Good spatial papers say which job geography is doing before estimating.
\end{frame}

\begin{frame}{Spatial clustering and Conley-style inference}
\small
\[
\widehat{\Var}(\hat\beta)
=
(X'X)^{-1}
\left[
\sum_i \sum_j K\!\left(\frac{d_{ij}}{c}\right)
x_i x_j' \hat u_i \hat u_j
\right]
(X'X)^{-1}
\]

\begin{itemize}
    \item Nearby residuals may share weather, labor demand, policy, infrastructure, or demographics.
    \item Choose coordinates, distance metric, cutoff \(c\), and kernel transparently.
    \item Report sensitivity to plausible spatial cutoffs.
    \item This fixes a variance problem, not bias from spillovers or sorting.
\end{itemize}
\end{frame}

\begin{frame}{Spillovers and exposure mappings}
\small
\[
Y_i = Y_i(D_i,E_i)
\qquad
E_i=\sum_{j\neq i}w_{ij}D_j
\]

\begin{itemize}
    \item Spatial work often violates simple no-interference logic.
    \item \(w_{ij}\) can encode distance, adjacency, commuting, school links, or market exposure.
    \item Miguel--Kremer is the canonical lesson: externalities are a causal object, not just noise.
    \item The mapping defines the estimand; it still needs a credible assignment or exposure design.
\end{itemize}
\end{frame}

\begin{frame}{Border designs and spatial RD}
\small
\[
Y_i =
\alpha+\tau\1\{b_i\ge 0\}
+f_-(b_i)\1\{b_i<0\}
+f_+(b_i)\1\{b_i\ge0\}
+X_i'\gamma+\varepsilon_i
\]

\begin{itemize}
    \item \(b_i\): signed distance to the policy or institutional boundary.
    \item Black: school-boundary comparisons and local amenities.
    \item Dell: historical boundary and continuity of predetermined geography.
    \item Dube--Lester--Reich: contiguous counties and local labor-market comparisons.
    \item Main caveat: borders bundle institutions, sorting, spillovers, and local shocks.
\end{itemize}
\end{frame}

\begin{frame}{Shift-share and place-based designs}
\small
\[
Z_\ell=\sum_s w_{\ell s}g_s
\]

\begin{itemize}
    \item \(w_{\ell s}\): baseline local shares; \(g_s\): common sector, origin, product, or task shock.
    \item The geography \(\ell\) defines the local exposure object.
    \item Diagnostics: dominant shares, effective shock variation, baseline year, leave-one-out shocks.
    \item Place-based policies require both treatment geography and exposure geography.
    \item Main caveat: polished exposure is not exogeneity.
\end{itemize}
\end{frame}

\begin{frame}{Sorting, MAUP, and correlated shocks}
\begin{itemize}
    \item When place is the treatment, selection into place is usually the central threat.
    \item Chetty--Hendren--Katz show why random assignment, move timing, and exposure duration matter.
    \item MAUP changes the estimand when tract, county, and commuting-zone estimates differ.
    \item Correlated local shocks can contaminate nearby controls or create spatial residual patterns.
    \item Geography sensitivity should be interpreted through the mechanism, not as a ritual.
\end{itemize}
\end{frame}

\begin{frame}{Design diagnostics checklist}
\small
\begin{enumerate}
    \item State the mechanism before choosing the geography.
    \item Define observation, treatment, exposure, and identifying variation separately.
    \item Write down the interference assumption.
    \item Check border balance, exposure support, or share decomposition.
    \item Diagnose contaminated controls and correlated local shocks.
    \item Report spatial inference as inference, not identification.
    \item Explain the counterfactual: nearby places, untreated time, other boundaries, or lower exposure.
\end{enumerate}
\end{frame}

\begin{frame}{Research Lab design}
\begin{columns}[T,onlytextwidth]
\begin{column}{0.32\linewidth}
\textbf{Reproduce}
\begin{itemize}
    \item synthetic border counties
    \item pair fixed effects
    \item treatment estimate
    \item robust vs Conley SE
\end{itemize}
\end{column}
\begin{column}{0.32\linewidth}
\textbf{Diagnose}
\begin{itemize}
    \item balance near borders
    \item cutoff sensitivity
    \item neighbor exposure
    \item role of geography
\end{itemize}
\end{column}
\begin{column}{0.32\linewidth}
\textbf{Transfer}
\begin{itemize}
    \item shift-share exposure
    \item dominant shares
    \item leave-one-out shocks
    \item design memo
\end{itemize}
\end{column}
\end{columns}
\end{frame}

\begin{frame}{What not to confuse}
\begin{itemize}
    \item Spatially robust standard errors are not a substitute for a spillover model.
    \item A border map is not proof of local continuity.
    \item A narrow bandwidth is not proof that sorting disappears.
    \item A shift-share index is not automatically an instrument.
    \item A stable estimate across geographies is reassuring only if each geography is meaningful.
\end{itemize}
\end{frame}

\begin{frame}{Takeaways}
\begin{itemize}
    \item Spatial causal inference is ordinary identification logic under harder geography.
    \item Place definitions, interference structure, and estimand choice are part of identification.
    \item Conley-style inference solves a dependence problem in uncertainty, not a design problem in bias.
    \item Border, exposure, shift-share, and mover designs each require paper-specific diagnostics.
    \item Lecture 17 asks when spatial equilibrium makes partial-equilibrium treatment effects incomplete.
\end{itemize}
\end{frame}

\end{document}
